% theme: aqueous_solutions
\MajorQuestion{水溶液}
\SetRowSpaceQanda

\Instruction{溶質、溶媒、溶液に関する次の問いに答えなさい。}
\QQ{物質が液体に溶けることを何というか。}{}
\QQ{液体に溶けている物質を\NamedBlank{solute}という。}{}
\QQ{\RefBlank{solute}を溶かしている液体を\NamedBlank{solvent}という。}{}
\QQ{\RefBlank{solute}が\RefBlank{solvent}に溶けた液体を\NamedBlank{solution}という。}{}
\QQ{\RefBlank{solvent}が水である\RefBlank{solution}を何というか。}{}

\Instruction{水溶液の性質と濃度に関する次の問いに答えなさい。}
\QQ{物質が溶けた水溶液は、見た目にはどのような状態になるか。}{}
\QQ{一つの水溶液では、どの部分を取っても濃さはどうなっているか。}{}
\QQ{水溶液全体の質量に対する\RefBlank{solute}の質量の割合を百分率で表したものを何というか。}{}
\QQ{質量パーセント濃度を求める式を、溶質と溶液の質量を用いて表しなさい。}{}
\QQ{\RefBlank{solution}の質量は、\RefBlank{solute}と\RefBlank{solvent}の質量からどのように求めるか。}{}

\Instruction{純粋な物質、混合物、ろ過に関する次の問いに答えなさい。}
\QQ{1種類の物質だけでできている物質を\NamedBlank{pure_substance}という。}{}
\QQ{複数の物質が混ざり合ってできたものを\NamedBlank{mixture}という。}{}
\QQ{液体と、それに溶けていない固体をろ紙で分ける操作を\NamedBlank{filtration}という。}{}
\QQ{\RefBlank{filtration}によって、ろ紙を通り抜けた液体を何というか。}{}
\QQ{ろ過するとき、液体を伝わせて静かに注ぐために使う器具は何か。}{}

\Instruction{飽和水溶液、溶解度、結晶に関する次の問いに答えなさい。}
\QQ{一定量の水に物質がそれ以上溶けない状態の水溶液を\NamedBlank{saturated_solution}という。}{}
\QQ{100\,gの水に溶ける物質の最大の質量を、その物質の\NamedBlank{solubility}という。}{}
\QQ{水の温度と\RefBlank{solubility}の関係を表したグラフを何というか。}{}
\QQ{純粋な物質で、いくつかの平面で囲まれた規則正しい形の固体を\NamedBlank{crystal}という。}{}
\QQ{固体を水に溶かし、再び\RefBlank{crystal}として取り出す操作を何というか。}{}
